
\medskip
\phantomsection\label{cor:sga3-VIB-1.3.1}%
\noindent\textbf{Corollary 1.3.1.---}
\emph{Let \(A\) be an Artinian local ring with residue field \(k\),
and let \(G\) be a flat \(A\)-group. The following assertions are
equivalent:}%
\footnote{\label{note:sga3-VIB-8}Editor's note: it is not enough to
suppose that \(G\) is flat over \(A\) at one point. For example,
suppose that \(A\ne k\), let \(H\) be the constant \(A\)-group
\((\mathbb Z/2\mathbb Z)_A\), and let \(G\) be the closed
\(A\)-subgroup of \(H\) whose nonidentity component is reduced.
Then the structural morphism \(G\to\operatorname{Spec}A\) is a
local isomorphism at the identity point \(e\), but is not flat at
the point \(g\ne e\).}
\begin{enumerate}
\item[\textup{\SGAInlineRefTarget{cor:sga3-VIB-1.3.1:item:i}(i)}] \emph{\(G\) is locally quasi-finite
  (respectively, unramified, smooth, or \'etale) over \(A\) at
  some point;}
\item[\textup{\SGAInlineRefTarget{cor:sga3-VIB-1.3.1:item:ii}(ii)}] \emph{\(G\) is locally quasi-finite
  (respectively, unramified, smooth, or \'etale) over \(A\).}
\end{enumerate}

\footnote{\label{note:sga3-VIB-9}Editor's note: we have added the
sentence that follows and simplified the proof of
\hyperref[lem:sga3-VIB-1.3.1.1]{Lemma~1.3.1.1} by invoking EGA~IV, 2.7.1 instead of
loc.\ cit., 17.7.5.}
\emph{Proof.}
If \(G\) satisfies one of the conditions in \textup{\SGARefLink{cor:sga3-VIB-1.3.1:item:i}{(i)}} at a
point \(x\), there is an open neighborhood \(U\) of \(x\) that is
of finite type over \(A\). Consequently, it is enough to apply
\hyperref[prop:sga3-VIB-1.3]{1.3} in the case where \(H\) is the unit \(A\)-group, taking the
following lemma into account.

\medskip
\phantomsection\label{lem:sga3-VIB-1.3.1.1}%
\noindent\textbf{Lemma 1.3.1.1.---}
\emph{Let \(A\) be an Artinian local ring and \(G\) an
\(A\)-group. If \(G\) contains a nonempty open subset of finite
type over \(A\), then \(G\) is locally of finite type over \(A\).}

\emph{Proof.}
By \hyperref[lem:sga3-VIB-1.3.0]{Lemma~1.3.0}, we may suppose that \(A\) equals its residue field
\(k\). Moreover, by \(\mathrm{fpqc}\) descent, we may suppose that
\(k\) is algebraically closed (cf.\ EGA~IV\(_2\), 2.7.1). Let
\(V\) be the open subset of \(G\) consisting of the points at
which \(G\) is of finite type over \(k\); by hypothesis,
\(V\ne\varnothing\). Since \(G\) is a Jacobson scheme, \(V\)
contains a closed point \(x\), and to prove that \(V=G\) it is
enough to show that every closed point \(y\) of \(G\) belongs to
\(V\). For such a point \(y\), the translation
\[
  r_y\circ r_x^{-1}
\]
sends \(x\) to \(y\), and hence \(y\in V\).

\medskip
\phantomsection\label{cor:sga3-VIB-1.3.2}%
\noindent\textbf{Corollary 1.3.2.---}
\emph{Let \(A\) be an Artinian local ring, and let \(u:G\to H\)
be a morphism between \(A\)-groups locally of finite type. The
following assertions are equivalent:
\begin{enumerate}
\item[\textup{\SGAInlineRefTarget{cor:sga3-VIB-1.3.2:item:i}(i)}] \(u\) is universally open;
\item[\textup{\SGAInlineRefTarget{cor:sga3-VIB-1.3.2:item:ii}(ii)}] \(u\) is open;
\item[\textup{\SGAInlineRefTarget{cor:sga3-VIB-1.3.2:item:iii}(iii)}] \(u\) is open at some point of \(G\);
\item[\textup{\SGAInlineRefTarget{cor:sga3-VIB-1.3.2:item:iv}(iv)}] the morphism
  \(u^0:G^0\to H^0\) induced by \(u\) is dominant;
\item[\textup{(iv\('\))}] \(u^0\) is surjective;
\item[\textup{\SGAInlineRefTarget{cor:sga3-VIB-1.3.2:item:v}(v)}] there is a connected component \(C\) of \(G\)
  such that, if \(D\) denotes the connected component of \(H\)
  containing \(u(C)\), the induced morphism
  \(u':C\to D\) is dominant.
\end{enumerate}}

\emph{Proof.}
The implications
\[
  \textup{(i)}\Longrightarrow\textup{(ii)}
  \Longrightarrow\textup{(iii)}
  \qquad\text{and}\qquad
  \textup{(iv\('\))}\Longrightarrow\textup{(iv)}
  \Longrightarrow\textup{(v)}
\]
are clear. Since \(G^0\) is of finite type over \(A\)
(Expos\'e~VI\(_A\), \SGARefLink{cor:sga3-VIB-2.4}{2.4}), and hence Noetherian, \(u^0\) is
quasi-compact. Thus \(u^0(G^0)\) is closed in \(H^0\) by \hyperref[prop:sga3-VIB-1.2]{1.2}, and
therefore \textup{\SGARefLink{cor:sga3-VIB-1.3.2:item:iv}{(iv)}} implies \textup{(iv\('\))}. On the other
hand, \(G^0\) (respectively, \(C\)) is open in \(G\)
(Expos\'e~VI\(_A\), \SGARefLink{cor:sga3-VIB-2.3}{2.3}), while \(H^0\) (respectively, \(D\)) is
irreducible (Expos\'e~VI\(_A\), \SGARefLink{cor:sga3-VIA-2.4.1}{2.4.1}). Hence \textup{\SGARefLink{cor:sga3-VIB-1.3.2:item:ii}{(ii)}}
implies \textup{\SGARefLink{cor:sga3-VIB-1.3.2:item:iv}{(iv)}} (respectively, \textup{\SGARefLink{cor:sga3-VIB-1.3.2:item:iii}{(iii)}} implies
\textup{\SGARefLink{cor:sga3-VIB-1.3.2:item:v}{(v)}}). It remains to prove that \textup{\SGARefLink{cor:sga3-VIB-1.3.2:item:v}{(v)}} implies
\textup{\SGARefLink{cor:sga3-VIB-1.3.2:item:i}{(i)}}.

We endow the open subset \(C\) (respectively, \(D\)) of \(G\)
(respectively, \(H\)) with its induced scheme structure, and let
\(u'\) denote the induced morphism \(C\to D\). Let \(k\) be the
residue field of \(A\).%
\footnote{\label{note:sga3-VIB-10}Editor's note: we have expanded
the original in what follows.}
Since the base change
\[
  \operatorname{Spec}k\longrightarrow\operatorname{Spec}A
\]
is a universal homeomorphism, we may suppose that \(A=k\). By
hypothesis \(u'\) is dominant; and since \(C\) is of finite type
over \(k\) (Expos\'e~VI\(_A\), 2.4.1), and hence Noetherian,
\(u'\) is quasi-compact. By EGA~IV\(_2\), 2.3.7,
\[
  u'\otimes_k\bar{k}
\]
is again quasi-compact and dominant, where \(\bar{k}\) denotes an
algebraic closure of \(k\). Since
\(C\otimes_k\bar{k}\) is a union of connected components of
\(G\otimes_k\bar{k}\), the morphism
\[
  u\otimes_k\bar{k}:
  G\otimes_k\bar{k}\longrightarrow H\otimes_k\bar{k}
\]
satisfies assertion \textup{(v)}. In view of
EGA~IV\(_2\), 2.6.4, we are thus reduced to the case where
\(A=k\) is an algebraically closed field.

In this case we may further replace \(u\) by \(u_{\mathrm{red}}\),
and are thereby reduced to the case where \(H\) is reduced. Let
\(\xi\) (respectively, \(\eta\)) be the generic point of \(C\)
(respectively, \(D\)). Since \(u'\) is quasi-compact and dominant,
\[
  u'(\xi)=\eta
\]
(cf.\ EGA~IV\(_1\), 1.1.5). Since \(H\) is reduced, the local
ring \(\mathcal O_{H,\eta}\) is a field, and hence \(u'\) is flat
at \(\xi\).%
\footnote{\label{note:sga3-VIB-11}Editor's note: the original
invoked the generic flatness theorem (EGA~IV\(_2\), 6.9.1),
which is not necessary here.}
It follows from \SGARefLink{prop:sga3-VIB-1.3}{1.3} that \(u\) is flat. Moreover, since \(u\) is
locally of finite type and \(H\) is locally Noetherian, \(u\) is
locally of finite presentation. Therefore \(u\) is universally
open by EGA~IV\(_2\), 2.4.6.

